\subsection{Testing for Large Extra Dimensions with Neutrino Oscillations --- arXiv:1101.0003}

\textbf{Authors:} P. A. N. Machado, H. Nunokawa, R. Zukanovich Funchal

\paragraph{主な主張}
CHOOZ・KamLAND・MINOSの結合解析でLEDによる有意な適合改善はなく、最軽質量がゼロの極限で半径を約$1.0\,\mu\mathrm{m}$（正常階層）と$0.6\,\mu\mathrm{m}$（逆階層）以下に99\%信頼水準で制限する\cite{Machado:2011LEDTests}。

\paragraph{新規性}
通常の振動パラメータとKK塔の効果を同時に扱う実験解析により、振動極小の移動・高速振動・全体的な消失を半径と絶対質量の制約に結び付けた。

\paragraph{位置付け}
最小LED模型の現代的な振動解析の基礎であり、原子炉と長基線加速器のflavour・質量階層に対する感度の違いを理解する際に有用である。
